\scanpage{58}
% @ lines 1–7 / text
$S$ schéma. On va définir des équivalences entre les \deleted{\unsure{sept}}\inserted{neuf} catégories suivantes attachées à $S$ -- et pour $S$ variable, ce seront des équivalences mutuelles de catégories fibrées sur $\mathrm{Sch}_{/S}$ (en fait, des gerbes \unsure{fpqc|fppf|fppc} sur $\mathrm{Sch}_{/S}$), formant un ``syst. transitif d'équivalences à isom can près'' :
% @ lines 8–26, two-column dictionary and return arrow / diagram
\diagram{\begin{tikzcd}[ampersand replacement=\&,column sep=0.9em,row sep=0.5em,cells={anchor=west}]
% item line 8: torsors
 {C_{0_S}=\operatorname{Tors}(\mathrm{GP}(1)_S)} \&
 {\parbox{7cm}{catégorie des torseurs \inserted{à droite} $\unsure{P|R}$ sous $\mathrm{GP}(1)_S$}}\\
% item lines 9–12: forms of groups, three rows
 {\left\{\begin{array}{l}
 C_{1_S}=\operatorname{For}(\mathrm{GP}(1)_S)\\[1.3em]
 C'_{1_S}=\operatorname{For}(\mathrm{Sl}(2)_S)\\[1em]
 \deleted{C''_{1_S}=\operatorname{For}(\mathrm{Gl}(2)_S)}
 \end{array}\right.} \&
 {\begin{array}{l}
 \parbox{7cm}{catégorie des ``formes'' $G$ (\unsure{fpqc|fppf} ou \unsure{f.ét.|f(ét)}) du \deleted{$\mathrm{GP}(1)$} schéma en \unsure{groupe|groupes} $\mathrm{GP}(1)$}\\[.5em]
 \text{catégorie des formes }\widetilde G\text{ de }\mathrm{Sl}(2)_S\\[1em]
 \deleted{\text{catégorie des formes }\unsure{\bar G|G}\text{ de }\mathrm{Gl}(2)_S}
 \end{array}}\\
% item lines 13–16: Lie algebras, stars at left
 {\begin{array}{r} \unsure{c}(*)\\[1em](*)\\[.8em]\deleted{(*)}\end{array}
 \left\{\begin{array}{l}
 C_{2_S}=\operatorname{For}(\mathfrak{gp}(1)_S)\\[1em]
 C'_{2_S}=\operatorname{For}(\mathfrak{sl}(2)_S)\\[.8em]
 \deleted{C''_{2_S}=\operatorname{For}(\mathfrak{gl}(2)_S)}
 \end{array}\right.} \&
 {\begin{array}{l}
 \parbox{7cm}{catégorie des formes $\unsure{\mathfrak g|g}$ de l'Algèbre de Lie $\mathfrak{gp}(1)_S$ de $\mathrm{GP}(1)_S$}\\
 \parbox{7cm}{\deleted{catégorie des formes} \rule[0.3ex]{1.4em}{0.3pt}\ \rule[0.3ex]{1.4em}{0.3pt}\ \rule[0.3ex]{1.4em}{0.3pt}\ $\widetilde{\mathfrak g}$ de $\mathfrak{sl}(2)_S\simeq\underline{\operatorname{Lie}}\mathrm{Sl}(2)_S$}\\
 \deleted{\text{catégorie des formes de }\mathfrak{gl}(2)_S=\underline{\operatorname{Lie}}\mathrm{Gl}(2)_S}
 \end{array}}\\
% item lines 17–19: quaternion algebras
 {C_{3_S}=\deleted{\unsure{\operatorname{Quat}(S)}}\operatorname{Quat}(\mathcal O_S)} \&
 {\parbox{7cm}{catégorie des Algèbres d'Azumaya\inserted{$M$} de rang $4$ sur $S$ (ou Algèbres de Quaternions)}}\\
% item lines 19–23: projective lines and quadrics
 {\left\{\begin{array}{l}
 C_{5_S}=\operatorname{Drpr}(S)\\[2.6em]
 C'_{5_S}=\operatorname{Quad}(S)
 \end{array}\right.} \&
 {\begin{array}{l}
 \parbox{7cm}{catégorie des fibrés en droites projectives $X$ sur $S$}\\[1em]
 \parbox{7cm}{catégorie des couples $(P,X)$ d'un fibré en plans projectifs sur $S$ et d'une quadrique lisse $X\subset P$}
 \end{array}}\\
% item lines 24–26: special orthogonal modules
 {\left[C_{4_S}=\overset{\operatorname{Modspquad}(S)^3}{\deleted{\unsure{\operatorname{Spéc}(S)_3}}}\right.}
 \arrow[uu,bend left=44,start anchor=west,end anchor=west] \&
 {\parbox{7cm}{catégories des ``Modules spéciaux orthogonaux de rang $3$ sur $S$'' $(E,q,\omega)$, $E$ fibré vect. de rg $=3$, $\omega$ \deleted{\unsure{forme}} base de $\det(\underline E)$, $q$ forme \unsure{pol$^{\mathrm{le}}$|quad.}, $\unsure{d(q)|\delta(q)}=-\omega^{\otimes2}$}}
\end{tikzcd}}
% archive sheet number: 57
